% v005 stage C draft (2026-11-17), narrative "empirical study" (author's decision 2026-11-16). Numbers from
% notes/analysis-v005/reports/01 (tests), 02 (stage B), 03 (RNG), 04 (E1), all computed from
% results/runs/c674c91e{,-vizier,-e1}/analysis/. Tables: _generated_*.tex (scripts/render_v005_tex.py,
% scripts/analyze_e1_v005.py). Retraining (K10): reports/06, results/retraining/.

\section{Results}
\label{sec:results}

All tests are those of the pre-registration: two-sided Mann--Whitney $U$ per (space, budget) cell, Holm--Bonferroni
within the cell's family, Cliff's $\delta$ oriented so that positive favours P3Net, 30 seeds per arm. Counts of
significant comparisons across cells are descriptive. The four primary spaces are JAHS-Bench-201 on CIFAR-10 (C10),
Colorectal Histology (CRC) and Fashion-MNIST (FMN), measured by normalised hypervolume, and NAS-HPO-Bench-II (NHB),
measured by IGD+ against its exact Pareto front.

\paragraph{Where P3 stands (Q1).} Table~\ref{tab:v005-q1-jahs-c10} and its companions give $\delta$ of P3Net against
the 17 comparison arms. The pattern is set by the budget more than by the space (Figure~\ref{fig:v005-rank-budget}). At $b = 50$ and $100$, P3Net is worse
than 12--17 of the 17 arms on each JAHS space, random search included ($\delta$ between $-0.43$ and $-0.68$), and
better than none; on NHB nearly all comparisons are not significant. At $b = 200$ and $350$ the picture reverses on
CRC and FMN: P3Net is worse than no arm and better than 2--8 of the 14, and its median ranks first or second of 19
arms in three of these four cells. C10 stays harder: at $b = 350$ P3Net is better than one arm and worse than four.
On NHB at $b = 350$ it is better than three and worse than five, with population-based evolutionary methods
(NSGANetV2, NSGA-III, SH-EMOA) ahead. Over all 16 cells, the arms most often better than P3Net are MO-BANANAS,
NSGA-Net, NSGA-III and SH-EMOA (9 cells each) and NSGANetV2 (8). The front P3Net returns is strong at its accurate end
and weaker at its cheap end: in the hypervolume decomposition on JAHS at $b = 350$ it ranks 7th, 3rd and 3rd of 24
arms on the accurate end of C10, CRC and FMN, against 15th, 11th and 4th on the cheap end
(Figure~\ref{fig:v005-fronts} shows the fronts of median runs in one such cell). The Gaussian-process
arms, run only up to $b = 100$, are among the strongest arms at those budgets: OSS Vizier has the best median of the 22 arms in 11 of
the 18 (space, budget) cells of all nine spaces (Table~\ref{tab:v005-s-gp}).

\paragraph{What the surrogate contributes (Q2, Q3).} No comparison of P3Net with canonical P3, P3Net-gate or
P3Net-nolink is significant (0 of 48; Table~\ref{tab:v005-q2-q3}), and P3Net is equivalent to its two variants at
the margin $m = 0.33$ in all 16 cells. This equivalence is largely mechanical. P3Net spends its budget on the
first-improvement hill climber and reaches the mixing stage late: the median number of mixing steps per run on the
primary spaces is 0, 0, 1 and 8 at $b = 50$, 100, 200 and 350, so P3Net-gate and P3Net-nolink end with the same final
value as P3Net in 94.6\,\% and 96.0\,\% of runs. With no mixing step P3Net \emph{is} canonical P3: the surrogate
takes one draw from the random generator when it is built, which shifts the random stream, and with that offset
removed the two produce identical histories (6 of 6 runs checked; without it, canonical P3 diverges in 5 of
the 6). That P3Net and canonical P3 end identically in only
5.8\,\% of runs therefore reflects mostly their different random streams, not surrogate decisions, and at these
budgets the Q1 results describe canonical P3 as much as P3Net.

\paragraph{The surrogate with room to act (E1, exploratory).} To see whether the surrogate matters once it acts, an
exploratory addendum ran the four P3 variants and random search at $b = 1000$ on new seeds
(Table~\ref{tab:v005-e1}). There P3Net takes 91--228 mixing steps per run, of which the surrogate screens 86--92\,\%,
and runs are no longer copies of each other (3--17\,\% identical to canonical P3). Still, none of the 12
comparisons is significant; at $m = 0.33$ nine are equivalent, two inconclusive, and one---canonical P3 on C10,
$\delta = -0.26$---has a confidence interval excluding zero in canonical P3's favour, while its Mann--Whitney test is
not significant ($p = 0.084$). The surrogate raises P3Net's optimiser time from 0--2.5\,s to 50--85\,s per run. All
four variants beat random search ($\delta$ 0.92--1.00). Within the design tested, a pairwise regressor of fitness
differences, the surrogate does not improve P3 on these benchmarks.

\paragraph{Multi-fidelity methods (Q4).} With the budget counted in full-evaluation equivalents, P3Net is worse than
all five multi-fidelity arms at $b = 50$ and $100$ on the three JAHS spaces (Table~\ref{tab:v005-q4-mf}). At
$b = 350$ it is better than ASHA, Hyperband and MOSH-BANANAS on CRC and FMN, and on FMN also than SH-EMOA. On NHB most
comparisons are not significant. Overall P3Net is better in 12, worse in 33 and indistinguishable in 35 of 80
comparisons.

\paragraph{Spaces outside the design loop (Q5).} On the four FCNet tasks, which played no part in designing P3Net,
the balance is more favourable than on the primary spaces: P3Net is better in 64, worse in 36 and indistinguishable in
148 of 248 comparisons (Table~\ref{tab:v005-q5-fcnet}). Protein structure is the exception, with 8 better against 21
worse. The primary-space result therefore does not transfer as a single verdict; it holds per space and per budget.

\paragraph{Optimiser cost (Q6).} Table~\ref{tab:v005-q6-costs} separates the optimiser's own time from the
benchmark queries. The arms strongest at small budgets are also the most expensive
(Figure~\ref{fig:v005-quality-cost}): at $b = 100$ on C10, OSS Vizier
needs a median of 1\,160\,s per run and 14\,GB of memory (up to 25\,GB elsewhere), NSGANetV2 624\,s, and qNEHVI 160\,s, against 0.3\,s for
P3Net; at $b = 350$ NSGANetV2 needs 2\,825\,s against 5.7\,s. On a tabular benchmark this cost dominates the run; in
real NAS, where one evaluation is a training run, it does not---the ratio matters when evaluations are cheap.
Measured in evaluations instead (Table~\ref{tab:v005-s-ctt}), P3Net is slow to reach a high target at $b = 350$: on
JAHS it reaches 99\,\% of the best median final HV in 67, 87 and 90\,\% of runs on C10, CRC and FMN, after a
median of 124, 135 and 127 evaluations, whereas most comparison arms reach it in all or nearly all runs and sooner
(on FMN after about 20). On NHB the target is hard for every arm: 95\,\% of the best median IGD+ is reached by at
most half of the runs of any arm (NSGANetV2) and by 3\,\% of P3Net's.

\paragraph{NAS-Bench-201 (descriptive).} With energy as the second objective, NAS-Bench-201 separates methods: median
IGD+ at $b = 350$ ranges from 0.0096 (NSGANetV2) to 0.0375 (MO-DSA), and the order of arms is stable between
$b = 200$ and 350 (Spearman 0.95). P3Net ranks 12th of 19 at both budgets, between canonical P3 and random search
(Table~\ref{tab:v005-s-nb201}).

\paragraph{P3 with eLyMPuS (Q7, supplement).} P3 with the eLyMPuS surrogate is worse than canonical P3 in 33 of 36
cells ($\delta$ down to $-0.38$; Table~\ref{tab:v005-s-q7}): at these budgets, discovering dependencies costs evaluations that canonical P3 spends
on search.

\paragraph{Secondary analyses (supplement).} The pre-registered descriptive analyses agree with the tests.
The anytime curves (Figure~\ref{fig:v005-s-anytime}) show why P3 trails at small budgets: on JAHS, P3Net and canonical
P3 rise late and in steps, the first evaluations going into hill-climbing single solutions, and they join the leading
group only after about 110--160 evaluations on CRC and FMN, and on C10 end below the best arm. Simple regret of $f_1$ at $b = 350$
(Table~\ref{tab:v005-s-regret}) matches the strong accurate end: on CRC and FMN P3Net's median regret (0.29 and
0.15 percentage points) is within 0.01 of the lowest of any arm, while on C10 and NHB it lies inside the range of the
other arms. The gap between test and validation error of the final fronts (Table~\ref{tab:v005-s-gap}) varies far more
between spaces than between arms (CRC $+1.0$ to $+1.9$, NAS-Bench-201 $-2.4$ to $-2.9$ percentage points). Recomputing HV with the reference points
$(1.05, 1.05)$ and $(1.2, 1.2)$ leaves the order of arms practically unchanged (Spearman $\ge 0.992$ in all 12 JAHS
cells) and P3Net's rank unchanged (Table~\ref{tab:v005-s-hvsens}). Every arm exhausts its budget in every run except
REINFORCE (15.8\,\% of runs; its policy converges and re-proposes one configuration), MO-DSA and MO-PSA (80.9 and
88.0\,\%; 1\,000 iterations without a new evaluation), MOTPE (99.6\,\%) and SH-EMOA (MF), whose schedule ends just
short of the budget (Table~\ref{tab:v005-s-budget}).

\paragraph{Retraining the knee points.} To check that the tabular values carry over to real training, we retrained
from scratch, once each, the knee point of every arm's median run at its largest budget on NAS-Bench-201 and
NAS-HPO-Bench-II with the benchmarks' training protocols (40 trainings for 49 knee points, as arms with the same
knee share one). On NAS-Bench-201 the retrained validation accuracy is close to the benchmark's value: mean
difference $-0.10$ percentage points, median absolute difference 0.24, at most 0.56, and Spearman 0.75 between the
two over the 16 trainings. On NAS-HPO-Bench-II it is lower by 0.88 points on average (median absolute difference 0.96,
at most 3.4), with Spearman 0.92 over 24 trainings; our network for this benchmark is built with a different code
base from the authors' (pre-registration \S8). Many knees lie close together (16 of 22 within 0.6 points on
NAS-Bench-201; 13 of 27 within one point of P3Net's on NAS-HPO-Bench-II), so a single training per knee reorders them: P3Net's knee ranks 9th among the arms by the benchmark's value on both benchmarks and
15th of 22 and 16th of 27 after retraining. The retraining thus confirms the benchmark values to within about a
point, and does not separate the arms.
